\section{Optimization Passes}

A \migaki{} execution plan is transformed by passes over \mir, while observed runs are analyzed by passes over trajectory graphs. v0.4 separates these stages so the system can earn optimization claims incrementally:

\begin{table}[t]
\centering
\small
\begin{tabular}{p{0.24\linewidth}p{0.62\linewidth}}
\toprule
\textbf{Stage} & \textbf{Representative passes} \\
\midrule
Observe & trace ingestion, span classification, cache-key recording, usage capture, error capture, local graph/report writing \\
Compare & exact cache-key matching, dependency-hash comparison, runtime-hash comparison, repeated operation reporting, longest reusable path detection \\
Estimate & avoidable token estimates, avoidable model/tool call counts, cost estimates, latency estimates, failed or blocked reuse reasons \\
Replay & exact model-call replay, deterministic tool replay, validator-gated output reuse, privacy-aware cache reads, stale-output rejection \\
Optimize & context deduplication, prompt-cache layout, routing, fallback planning, retry planning, parallelism, provider-aware lowering \\
\bottomrule
\end{tabular}
\caption{Evidence-first optimization stages for \migaki{}.}
\end{table}

The earliest pass interface can remain small:

\begin{lstlisting}[language={}]
interface MigakiPass {
  name: string
  version: string
  apply(input: MIRPlan | MigakiGraph, context: PassContext):
    Promise<PassResult>
}

type PassResult = {
  output: MIRPlan | MigakiGraph | ComparisonReport
  diff?: MIRPlanDiff
  evidence: Evidence[]
  warnings: Warning[]
}
\end{lstlisting}

The important invariant is not the exact TypeScript shape. The invariant is that every pass emits evidence. A pass that changes execution without a diff, warning surface, validator boundary, or benchmarkable report is not an optimizer. It is hidden behavior.

\begin{principle}[First compare, then replay]
Graph comparison is a product milestone distinct from cache replay. \migaki{} should prove that reusable nodes can be detected before it proves that reusable nodes can be skipped.
\end{principle}
